\documentclass[11pt]{amsart}
\usepackage[T1]{fontenc}
\usepackage[utf8]{inputenc}
\usepackage{amssymb}
\usepackage[margin=1.05in]{geometry}
\usepackage[hidelinks]{hyperref}
\usepackage{booktabs}
\usepackage{longtable}
\usepackage{array}
\usepackage{datetime}
\numberwithin{equation}{section}
\newtheorem{theorem}{Theorem}[section]
\newtheorem{proposition}[theorem]{Proposition}
\newtheorem{lemma}[theorem]{Lemma}
\newtheorem{corollary}[theorem]{Corollary}
\theoremstyle{definition}
\newtheorem{definition}[theorem]{Definition}
\theoremstyle{remark}
\newtheorem{remark}[theorem]{Remark}
\DeclareMathOperator{\deh}{deh}
\DeclareMathOperator{\topc}{top}
\DeclareMathOperator{\ev}{ev}
\newcommand{\code}[1]{\texttt{#1}}
\newcommand{\ZZ}{\mathbb Z}
\AddToHook{shipout/foreground}{%
 \put(0.5\paperwidth,-\paperheight+0.55cm){\makebox[0pt]{\tiny\ttfamily
 unit/doc/\jobname.tex\ ---\ compiled \today\ \currenttime}}}
\title[Zariski's Main Theorem by finite section lifting]{A constructive proof of
Zariski's Main Theorem\\ by finite section lifting, and its formalisation in Unit}
\date{30 October 2026}
\begin{document}
\sloppy
\begin{abstract}
We prove the algebraic form of Zariski's Main Theorem used in the
local structure theory of \'etale algebras: if $E=R[x_1,\ldots,x_n]$ is
a finitely generated algebra over a commutative ring $R$, $I$ is an
ideal of $R$, and $E/IE$ is finite over $R/I$, then there is
$s\in1+IE$ such that $R[s,sx_1,\ldots,sx_n]$ is finite over $R$ and
has the same localisation at $s$ as $E$. The proof is constructive and
uses no hypothesis on $R$. It replaces the theorem on formal functions
by an explicit finite lifting argument in a \v Cech complex, driven by
constructive noetherianity of a polynomial ring over $\ZZ$. The last
section describes the complete machine-checked proof in Unit.
\end{abstract}
\maketitle


\section{The theorem}

All rings are commutative. ``Finite'' means finitely generated as a
module. ``There is'' means mere existence: the proof produces the
objects, but it does not claim that they are determined by the data.

\begin{theorem}\label{thm:main}
Let $R$ be a ring, $I\subseteq R$ an ideal, and $E$ an $R$-algebra
generated by $x_1,\ldots,x_n$. Suppose that $E/IE$ is finite as an
$R/I$-module. Then there is $s\in E$ such that
\begin{enumerate}
\item $s-1\in IE$;
\item $B=R[s,sx_1,\ldots,sx_n]\subseteq E$ is a finite $R$-module;
\item the map $B[1/s]\to E[1/s]$ induced by the inclusion is an
isomorphism.
\end{enumerate}
\end{theorem}

No noetherian, coherence, field or dimension hypothesis is imposed on
$R$, and $E$ need not be finitely presented. With $I=\mathfrak m$ for
a residually discrete local ring, every prime of $E$ above $\mathfrak m$
avoids $s$, which is the form used in the local structure theorem for
\'etale algebras.

\subsection*{Strategy}
Classically one embeds $\operatorname{Spec}E$ as the chart $Z\neq0$
of a projective scheme $W$, notes that $Z^{-1}$ is a section of
$\mathcal O_W(-1)$ over each infinitesimal neighbourhood $W\times R/I^k$
of the special fibre (the affine part is open and closed there), and
lifts these sections to one global section by the theorem on formal
functions. Clearing denominators in that section gives the result.
We keep the shape of this argument but replace every infinite object
by a finite certificate:
\begin{itemize}
\item the integrality conclusion follows from a few
\emph{degree-lowering relations} (Section~\ref{sec:certificate});
\item an arbitrary base reduces to a polynomial ring over $\ZZ$, which
is constructively noetherian (Section~\ref{sec:descent});
\item the sections $Z^{-1}$ modulo $I^k$ are given by explicit
formulas (Section~\ref{sec:sections});
\item the formal-functions step becomes one finite dependence relation,
produced by a noetherian argument over the Rees algebra
(Section~\ref{sec:lifting});
\item the corrected section yields the degree-lowering relations
(Section~\ref{sec:extraction}).
\end{itemize}
We present the proof backwards from its output.

\section{Certificates}\label{sec:certificate}

Let $A\to F$ be a ring map, $y_1,\ldots,y_n\in F$ and
$b_1,\ldots,b_t\in A$.

\begin{definition}\label{def:cert}
A \emph{ZMT certificate} for $(F,y,b)$ consists of $s\in F$ with
$s=1+\sum_j b_jv_j$ for some $v_j\in F$, together with monic
polynomials $Q_0,Q_1,\ldots,Q_n\in A[T]$ such that $Q_0(s)=0$ and
$Q_i(sy_i)=0$ for every $i$.
\end{definition}

\begin{lemma}\label{lem:cert-conclusion}
Suppose $E$ is generated by $x_1,\ldots,x_n$ over $R$ and $(E,x,a)$
has a certificate with all $a_j\in I$. Then $s$ satisfies the
conclusion of Theorem~\ref{thm:main}.
\end{lemma}
\begin{proof}
Clearly $s-1\in IE$. Put $w_0=s$, $w_i=sx_i$. Each $w_j$ satisfies
the monic equation $Q_j$, of degree $d_j$; hence the monomials
$w^e$ with $e_j<d_j$ span $B$ over $R$. The map
$B[1/s]\to E[1/s]$ is injective because $B\to E$ is injective. Every
$e\in E$ is $H(x)$ for a polynomial $H$ of some degree $N$, and then
$s^NH(x)$ is a polynomial in the $w_i$, so $e/s^k=(s^Ne)/s^{N+k}$ is
in the image.
\end{proof}

\begin{lemma}[Determinant trick]\label{lem:trick}
Let $V\subseteq F$ be the $A$-span of a finite family containing $1$.
If $u\in F$ satisfies $uV\subseteq V$, then $u$ is a root of a monic
polynomial over $A$.
\end{lemma}
\begin{proof}
Write $uv_k=\sum_l\alpha_{lk}v_l$. By the adjugate identity, the
characteristic polynomial $\chi$ of $(\alpha_{lk})$ satisfies
$\chi(u)v_k=0$ for all $k$, in particular $\chi(u)=\chi(u)\cdot1=0$.
The family need not be independent.
\end{proof}

Let $\deh:A[X_1,\ldots,X_n,Z]\to F$ be the ring map $Z\mapsto1$,
$X_i\mapsto y_i$ (``dehomogenisation'').

\begin{proposition}[Degree lowering]\label{prop:dl}
Suppose $s=1+\sum_jb_jv_j$, $s=\deh(\sigma)$ with $\sigma$ homogeneous
of degree $D$, and for every $i$
\begin{equation}\tag{DL}
 s\,y_i^{a_i+1}=\deh(q_i),\qquad q_i \text{ homogeneous of degree } a_i .
\end{equation}
Then $(F,y,b)$ has a certificate with this $s$.
\end{proposition}
\begin{proof}
Let $d=D+\sum_ia_i$ and let $V$ be the span of the monomials $y^\beta$
with $|\beta|\le d$; it contains $1$. It suffices to show
$s\,y^\beta\in V$ whenever $|\beta|\le d+1$, since then $sV\subseteq V$
and $(sy_i)V\subseteq V$, and Lemma~\ref{lem:trick} applies. Either
$\beta_i\ge a_i+1$ for some $i$, and then $y^\beta=y^\gamma y_i^{a_i+1}$
and $sy^\beta=y^\gamma\deh(q_i)$ has degree at most
$|\gamma|+a_i=|\beta|-1\le d$; or $|\beta|\le\sum_ia_i$, and then
$sy^\beta=\deh(\sigma)y^\beta$ has degree at most $D+\sum_ia_i=d$.
This dichotomy is decided on the exponents only.
\end{proof}

Only one finite module is used, for $s$ and all $sy_i$ at once. No
ring of global sections and no multiplication table of $B$ is
computed.

\section{Reduction to residual data over \texorpdfstring{$\ZZ[C]$}{Z[C]}}\label{sec:descent}

\begin{definition}
Residual data over a ring $A$ in $n$ variables consist of
$b_1,\ldots,b_t\in A$, monic $P_1,\ldots,P_n\in A[T]$ and polynomials
$H_{ij}\in A[Y_1,\ldots,Y_n]$. The residual algebra is
\[
 F(\mathcal D)=A[Y_1,\ldots,Y_n]\big/\bigl(P_i(Y_i)-\textstyle\sum_jb_jH_{ij}(Y)\bigr)_{i},
\]
with generators $y_i$, the classes of $Y_i$.
\end{definition}

\begin{lemma}\label{lem:residual}
Under the hypotheses of Theorem~\ref{thm:main} there are residual
data over $R$ with all $b_j\in I$ and $P_i(x_i)=\sum_jb_jH_{ij}(x)$
in $E$.
\end{lemma}
\begin{proof}
Lift finitely many generators of $E/IE$ to $g_1,\ldots,g_N\in E$ and
write $1\equiv\sum_k\alpha_kg_k$ and $x_ig_k\equiv\sum_l\alpha^{(i)}_{lk}g_l$
modulo $IE$. The characteristic polynomial $P_i$ of
$(\alpha^{(i)}_{lk})$ gives $P_i(x_i)g_k\in IE$ for all $k$, hence
$P_i(x_i)\in IE$. Collect the finitely many coefficients from $I$ in
these memberships into one list $b$, and write the other factors as
polynomials in $x$, using that $x$ generates $E$.
\end{proof}

Only finitely many lifts and polynomial expressions are chosen.

\begin{definition}
Let $A_0=\ZZ[C_1,\ldots,C_M]$ where the $C_k$ stand for all
coefficients occurring in residual data $\mathcal D$ over $R$: the
$b_j$, the non-leading coefficients of the $P_i$ and the coefficients
of the $H_{ij}$, each at its own index. Replacing each coefficient by its
variable gives universal residual data $\mathcal D_0$ over $A_0$
(leading coefficients stay $1$), and $h:A_0\to R$, $C_k\mapsto c_k$,
maps $\mathcal D_0$ to $\mathcal D$.
\end{definition}

The residual equations of $\mathcal D_0$ hold for $x$ in $E$ through
$h$, so there is $\varphi:F(\mathcal D_0)\to E$ with $\varphi(y_i)=x_i$.
Neither $h$ nor $\varphi$ need be injective, and no kernel is computed.

\begin{lemma}[Transport]
A certificate for $(F(\mathcal D_0),y,b_0)$ over $A_0$ is sent by
$(h,\varphi)$ to a certificate for $(E,x,b)$ over $R$.
\end{lemma}
\begin{proof}
$\varphi(1+\sum b_{0j}v_j)=1+\sum_jb_j\varphi(v_j)$ and
$(hQ)(\varphi(u))=\varphi(Q(u))$; monicity is preserved.
\end{proof}

Recall that a ring is CFN if it is coherent, strongly discrete and
constructively noetherian, in the sense of Coquand and Persson. The
ring $\ZZ$ is CFN, and so is a polynomial ring over a CFN ring
(Hilbert's basis theorem in this form). Hence $A_0$ is CFN. By
Lemmas~\ref{lem:cert-conclusion}, \ref{lem:residual} and transport,
Theorem~\ref{thm:main} reduces to:

\begin{theorem}[Universal ZMT]\label{thm:universal}
For every CFN ring $A$, every $n$ and all residual data $\mathcal D$
over $A$ in $n$ variables, $(F(\mathcal D),y,b)$ merely has a
certificate.
\end{theorem}

\begin{lemma}[Padding]\label{lem:padding}
Theorem~\ref{thm:universal} for $n+1$ variables implies it for $n$
variables.
\end{lemma}
\begin{proof}
Add a variable $Y_{n+1}$ with $P_{n+1}=T$ and $H_{n+1,j}=0$. The new
residual algebra maps to $F(\mathcal D)$ by $Y_{n+1}\mapsto0$; transport a
certificate and forget $Q_{n+1}$.
\end{proof}

So it suffices to treat $n=m\ge2$, which gives at least three
homogeneous coordinates below. For the rest of the proof fix a CFN ring
$A$, residual data $\mathcal D$ in $m\ge2$ variables, $F=F(\mathcal D)$,
and $I=(b_1,\ldots,b_t)\subseteq A$.

\section{The projective setting}\label{sec:projective}

\subsection{Homogeneous equations}
Use homogeneous coordinates $X_1,\ldots,X_m,Z$ and
$S=A[X_1,\ldots,X_m,Z]$. Let $d_i=\deg P_i$, choose $\delta_i$ larger
than the degrees of all $H_{ij}$, and let $e_i=d_i+\delta_i$. Put
\[
 M_i=Z^{d_i}P_i(X_i/Z)=X_i^{d_i}+Z\bar Q_i,\qquad
 \mathbf P_i=Z^{\delta_i}M_i-\sum_jb_j\,Z^{e_i}H_{ij}(X/Z),
\]
homogeneous of degree $e_i$, and $J=(\mathbf P_1,\ldots,\mathbf P_m)$.
Then $\deh(\mathbf P_i)=0$ in $F$: these are exactly the residual
equations. We do not saturate $J$; the projective scheme it defines may
have extra components at infinity.

\subsection{\v Cech cochains}
For a nonempty set $\sigma$ of coordinates let $S_\sigma$ be the
localisation of $S$ at the product of the coordinates in $\sigma$, and
$S_\sigma(k)$ its part of degree $k$. Then
\[
 C^p(k)=\prod_{|\sigma|=p+1}S_\sigma(k)
\]
with the alternating differential $d$ form the \v Cech complex of
$\mathcal O(k)$ on the standard charts. We represent an element of
$S_\sigma(k)$ by its finitely supported Laurent coefficient function.
This representation is faithful over every ring, including rings with
zero divisors, because multiplication by a monomial only translates
coefficients.

\begin{definition}
$JC^p(k)$ is the image of
$\bigoplus_iC^p(k-e_i)\to C^p(k)$, $(y_i)\mapsto\sum_i\mathbf P_iy_i$.
$I^rC^p(k)$ is the set of finite sums $\sum b^\alpha\omega_\alpha$ with
$|\alpha|=r$. We write $x\equiv y$ for $x-y\in JC^p(k)$.
\end{definition}

These are images inside the complex. They are not tensor products,
and no flatness is assumed or needed. Both are preserved by $d$.

The section $Z^{-1}$ appears through the cochain $z^{\mathrm{init}}$
with $z^{\mathrm{init}}_Z=Z^{-1}$. The main intermediate result is:

\begin{theorem}[Closed lift]\label{thm:closed}
There is $\zeta\in C^0(-1)$ with
\[
 d\zeta\in JC^1(-1),\qquad \zeta_Z-Z^{-1}\in IS_Z(-1)+JS_Z(-1).
\]
\end{theorem}

Sections~\ref{sec:sections} and~\ref{sec:lifting} prove it.
Section~\ref{sec:extraction} shows that it implies
Theorem~\ref{thm:universal}.

\section{Finite-fibre sections}\label{sec:sections}

\begin{proposition}\label{prop:sections}
For every $r\ge0$ there is an explicit $z_r\in C^0(-1)$ with
\begin{enumerate}
\item $(z_r)_Z=Z^{-1}$;
\item $dz_r\in I^{r+1}C^1(-1)+JC^1(-1)$;
\item $z_r-z_0\in IC^0(-1)+JC^0(-1)$.
\end{enumerate}
The witnesses for these memberships are explicit as well.
\end{proposition}

\begin{proof}
On the chart $X_i$ put $u=Z/X_i\in S_{X_i}(0)$. Dividing by
$X_i^{e_i}$,
\[
 \mathbf P_i/X_i^{e_i}=G_i-\sum_jb_jh_{ij},\qquad
 G_i=u^{\delta_i}H_i,\qquad H_i=M_i/X_i^{d_i}=1+uQ_i,
\]
with $Q_i,h_{ij}\in S_{X_i}(0)$. Hence $G_i\equiv\sum_jb_jh_{ij}$, and
\begin{equation}\label{eq:G}
 G_i^p\in I^p+J\qquad\text{for every }p.
\end{equation}
Fix $r$, let $p=r+1$, write $H_i^p=1+uQ_{i,p}$, and choose
$N\ge\delta_ip$ for all $i$. Define
\[
 t_i=(-Q_{i,p})^{N+1}u^N,\qquad (z_r)_{X_i}=t_i/X_i,\qquad (z_r)_Z=Z^{-1}.
\]
Then $u\,t_i=(1-H_i^p)^{N+1}=1-H_i^pa_i$ for a polynomial $a_i$ in
$H_i^p$, so
\begin{equation}\label{eq:unit}
 Z\,(z_r)_{X_i}=1-a_iH_i^p .
\end{equation}

\emph{Edge $\{Z,X_k\}$.} By \eqref{eq:unit},
$(z_r)_{X_k}-Z^{-1}=-a_kH_k^p/Z$. On $S_{ZX_k}$ the element $u_k$ is
invertible and $H_k^p=u_k^{-\delta_kp}G_k^p\in I^p+J$ by \eqref{eq:G}.

\emph{Edge $\{X_k,X_l\}$.} Write $z_k=(z_r)_{X_k}$. Since
$u_lt_l+a_lH_l^p=1$,
\[
 z_k=z_ku_lt_l+z_ka_lH_l^p,\qquad z_l=z_lu_kt_k+z_la_kH_k^p,
\]
and $z_ku_lt_l=t_kt_lZ/(X_kX_l)=z_lu_kt_k$. Hence
\[
 z_k-z_l=z_ka_lH_l^p-z_la_kH_k^p .
\]
In the first term $t_k$ contains $u_k^N=(X_l/X_k)^Nu_l^N$, and
$u_l^NH_l^p=u_l^{N-\delta_lp}G_l^p\in I^p+J$. The second term is
symmetric. This proves (2) with $p=r+1$.

\emph{Compatibility.} On $Z$ the difference $z_r-z_0$ vanishes. On
$X_k$ apply the same cross-term identity to $z_r$ and $z_0$: the
products $z_{r,k}\,u\,t_{0,k}$ and $z_{0,k}\,u\,t_{r,k}$ agree, and
each remaining term contains a factor $u^{N}H_k$ with $N\ge\delta_k$,
which lies in $I+J$ by \eqref{eq:G} with $p=1$.
\end{proof}

No idempotent, no comparison between different bounds $N$, and no
hypothesis on $A$ is used. Classically, $t_i/X_i$ is $Z^{-1}$ on the
part of the chart where the fibre lies and $0$ elsewhere; the formulas
make this precise modulo $I^{r+1}$ without forming the decomposition.

\section{Finite lifting}\label{sec:lifting}

Here the CFN hypothesis on $A$ is used. Classically this step expresses
finiteness of $\bigoplus_rH^1(W,I^r\mathcal O_W(-1))$ over the Rees
algebra. We need only one relation among the classes of the $dz_r$.

\subsection{The Rees presentation}
Let $D=A[U_1,\ldots,U_t]$, graded by $U$-degree, and
$\ev:D\to A$, $U_j\mapsto b_j$. Elements of degree $k$ evaluate into
$I^k$. Consider
\[
 \rho:D[X,Z]\longrightarrow(S/J)[T],\qquad U_j\longmapsto b_jT .
\]
Its kernel $K$ is bihomogeneous, for the $(X,Z)$-degree and the
$U$-degree. A polynomial $F$ of $U$-degree $k$ satisfies
$\rho(F)=F(b,X,Z)T^k$, so
\begin{equation}\label{eq:kernel}
 F\in K\iff F(b,X,Z)\in J .
\end{equation}

\begin{lemma}
$K$ is generated by a finite list $g_1,\ldots,g_a$ of bihomogeneous
polynomials, of degrees $\delta(g_i)$ in $(X,Z)$ and $f_i$ in $U$, with
explicit identities $g_i(b,X,Z)=\sum_q h_{iq}\mathbf P_q$.
\end{lemma}
\begin{proof}
$K=(J,U_1-b_1T,\ldots,U_t-b_tT)\cap A[U,X,Z]$ is an elimination ideal.
Since $A[U,X,Z]$ is CFN, eliminating the last variable $T$ gives
finitely many generators with membership witnesses. Replace each
generator by its bihomogeneous components, which lie in $K$ because
$K$ is bihomogeneous.
\end{proof}

We do not assume that $J$ and the relations $b_iU_j=b_jU_i$ generate
$K$: that would require flatness.

\begin{lemma}[Localised kernel membership]\label{lem:localised}
Let $c\in C^p_D(k)$ be homogeneous of $U$-degree $k'$ with
$\ev(c)\in JC^p_A(k)$. Then merely $c=\sum_ig_iu_i$ with
$u_i\in C^p_D(k-\delta(g_i))$ of $U$-degree $k'-f_i$.
\end{lemma}
\begin{proof}
Work on each of the finitely many charts $\sigma$. Clear one monomial
denominator $X^L$ common to $c_\sigma$ and to the $J$-multipliers of
$\ev(c_\sigma)$. The polynomial $X^Lc_\sigma$ has $U$-degree $k'$ and
lies in $J$ after evaluation, so it is in $K$ by \eqref{eq:kernel}.
Write it as $\sum g_iF_i$, take bihomogeneous components of the $F_i$,
and divide by $X^L$. A cochain is a family indexed by charts, so these
chart-wise lifts form a cochain; only finitely many choices are made.
\end{proof}

\subsection{Lifting the sections to the Rees algebra}
By Proposition~\ref{prop:sections}(2), $dz_r=w_r+j_r$ with $j_r\in JC^1$
and $w_r=\sum_{|\alpha|=r+1}b^\alpha\omega_\alpha$. Put
\[
 c_r=\sum_{|\alpha|=r+1}U^\alpha\omega_\alpha\in C^1_D(-1),
\]
of $U$-degree $r+1$, with $\ev(c_r)=w_r$. Then
$\ev(dc_r)=dw_r=-dj_r\in JC^2$, and Lemma~\ref{lem:localised} gives,
merely,
\begin{equation}\label{eq:lift}
 dc_r=\sum_ig_iu_{r,i}.
\end{equation}

\subsection{Noetherian dependence}

\begin{theorem}\label{thm:dependence}
Let $D$ be CFN, let $S_D=D[X_0,\ldots,X_{\nu-1}]$ have $\nu\ge3$
coordinates, and let $g_1,\ldots,g_a$ be homogeneous of degrees
$\delta_i\ge0$. Let $c_0,c_1,\ldots\in C^1_D(-1)$ each merely satisfy
$dc_r=\sum_ig_iu_{r,i}$. Then merely there are $N$ and
$q_0,\ldots,q_{N-1}\in D$ with
\[
 c_N-\sum_{i<N}q_ic_i\in d\,C^0_D(-1)+\sum_ig_i\,C^1_D(-1-\delta_i).
\]
\end{theorem}

The proof has four ingredients.

\emph{(a) Contraction of free modules.} For any ring $D$ and any twist
$k$, the complex $C^\bullet_D(k)$ is exact in degrees $0<p<\nu-1$, and in
the top degree $\nu-1$ a cochain is a coboundary when its coefficients
at the all-negative monomials vanish. The contraction is monomial by
monomial: a Laurent monomial whose set of negative coordinates is $\tau$
occurs only on charts $\sigma\supseteq\tau$. These charts form a cone,
hence a contractible complex, unless $\tau$ is empty or is everything.
In a fixed degree there are only finitely many all-negative monomials.

\emph{(b) A bounded homogeneous resolution.} Coherence of $S_D$ gives
homogeneous matrices
\[
 F_{\nu-2}\xrightarrow{A_{\nu-3}}\cdots\xrightarrow{A_1}F_1
 \xrightarrow{A_0=(g_1,\ldots,g_a)}F_0=S_D(-1)
\]
between finite sums of shifted free modules, each column list generating
the homogeneous kernel of the previous matrix. This is exact at
$F_1,\ldots,F_{\nu-3}$; nothing is claimed at the left end, and no
termination of a resolution is used. Exactness persists on every chart:
clear a common denominator, apply kernel generation, and translate back.

\emph{(c) Raising and reflection.} Given $dc=A_0u_1$ with
$u_1\in C^2(F_1)$, we have $A_0\,du_1=ddc=0$, so chart-wise exactness
gives $du_1=A_1u_2$. Continuing, we obtain $u_{\nu-2}\in C^{\nu-1}(F_{\nu-2})$,
in top degree. Let $\topc(u)\in D^v$ be its finitely many coefficients at
all-negative monomials. Define $\Gamma(c,p)$ to mean that such a chain
merely exists with $\topc(u_{\nu-2})=p$. Then $\Gamma$ is a
$D$-submodule of $C^1\times D^v$, and every $c_r$ has a partner. Moreover
\[
 \Gamma(c,0)\ \Longrightarrow\ c\in dC^0+A_0C^1(F_1).
\]
Indeed, if $\topc(u_{\nu-2})=0$, then (a) gives $u_{\nu-2}=dh_{\nu-2}$.
Then $u_{\nu-3}-A_{\nu-3}h_{\nu-2}$ is closed below the top degree, hence
equal to some $dh_{\nu-3}$. Descending this way, we reach
$u_1=A_1h_2+dh_1$. Then $c-A_0h_1$ is closed in degree $1$; since
$1<\nu-1$, it equals some $dh_0$. This is where $\nu\ge3$ is used.

\emph{(d) Dependence with partners under truncation.} Let $D$ be CFN
and $\Gamma\subseteq Z\times D^v$ a submodule such that $\Gamma(z,0)$
implies $z\in B$, for a submodule $B$. If every $b_r$ merely has a
partner, then merely $b_N-\sum_{i<N}q_ib_i\in B$ for some $N$ and $q_i$.
The proof is by induction on free lists in $D^v$, from constructive
noetherianity. At each step one mere partner is obtained inside the
propositional goal. If it lies in the span of the earlier partners
(decidable by strong discreteness), we have the relation; otherwise the
free list grows. No sequence of partners is ever chosen.

Theorem~\ref{thm:dependence} follows from (c) and (d), with
$Z=C^1_D(-1)$. It has no finite-matrix-kernel or cohomology-generator
computation: dependence is found directly in the ambient free module
$D^v$.

\begin{corollary}[Weighted and evaluated form]\label{cor:weighted}
For the $c_r$ of~\eqref{eq:lift} there are $N$ and
$r_i\in I^{N-i}$ $(i<N)$, $v\in I^{N+1}C^0(-1)$ and
\[
 w_N-\sum_{i<N}r_iw_i\in dv+JC^1(-1).
\]
\end{corollary}
\begin{proof}
Apply Theorem~\ref{thm:dependence} over $D=A[U]$, which is CFN, with
$\nu=m+1\ge3$. Since $c_r$ has $U$-degree $r+1$ and $g_i$ has $U$-degree
$f_i$, project the relation to $U$-degree $N+1$. Projection commutes
with $d$ and with multiplication by $U$-homogeneous polynomials, so
$q_i$ may be replaced by its part of degree $N-i$, the primitive by
its part of degree $N+1$, and the multipliers of $g_i$ by their parts of
degree $N+1-f_i$. Now evaluate at $U=b$. Here
$\ev(g_i)=\sum_qh_{iq}\mathbf P_q$, so the $g$-terms land in $JC^1$.
Negative degrees give zero parts.
\end{proof}

\subsection{Correction}
\begin{proof}[Proof of Theorem~\ref{thm:closed}]
Let $\zeta=z_N-\sum_{i<N}r_iz_i-v$. Then
\[
 d\zeta\equiv w_N-\sum_{i<N}r_iw_i-dv\equiv0
\]
modulo $JC^1(-1)$. Also $\zeta-z_0=(z_N-z_0)-\sum r_iz_i-v$ lies in
$IC^0+JC^0$, because $N-i\ge1$. On the chart $Z$ this gives
$\zeta_Z-Z^{-1}\in IS_Z+JS_Z$, since $(z_0)_Z=Z^{-1}$.
\end{proof}

\section{Extraction}\label{sec:extraction}

\begin{proposition}\label{prop:extraction}
A cochain $\zeta$ as in Theorem~\ref{thm:closed} gives degree-lowering
data (DL) for $(F,y,b)$.
\end{proposition}
\begin{proof}
Take one finite bound $\beta$ for the negative exponents occurring in
the supports of $\zeta$, of its $J$-multipliers on the edges
$\{Z,X_i\}$, and of the $I$- and $J$-multipliers in
$\zeta_Z-Z^{-1}$. Put $k_Z=\beta_Z+1$ and $k_i=\beta_i+1$. Then
\[
 \zeta_Z=\sigma/Z^{k_Z},\qquad (\zeta)_{X_i}=q_i/X_i^{k_i},
\]
with $\sigma$ and $q_i$ polynomials, homogeneous of degrees $\beta_Z$
and $\beta_i$. Put $s=\deh(\sigma)$.

On the edge $\{Z,X_i\}$ the condition $d\zeta\in JC^1(-1)$ reads
$q_i/X_i^{k_i}-\sigma/Z^{k_Z}=\sum_q\mathbf P_qm_q$. Multiplying by
$Z^{k_Z}X_i^{k_i}$ makes every term polynomial, by the choice of
$\beta$:
\[
 Z^{k_Z}q_i-X_i^{k_i}\sigma=\sum_q\mathbf P_q\,(Z^{k_Z}X_i^{k_i}m_q).
\]
Dehomogenising kills the right side and gives
$s\,y_i^{\beta_i+1}=\deh(q_i)$, which is (DL) with $a_i=\beta_i$.
Similarly, multiplying $\zeta_Z-Z^{-1}=\sum_jb_jv_j+\sum_q\mathbf P_qn_q$
by $Z^{k_Z}$ and dehomogenising gives $s=1+\sum_jb_j\deh(v'_j)$.
\end{proof}

Only the overlaps of the chart $Z$ with the other charts are used here.
All identities are proved as equalities of Laurent coefficients and
then evaluated, which is valid over every ring.

\section{Proof of Theorem~\ref{thm:main}}
For CFN $A$ and $m\ge2$, Proposition~\ref{prop:sections},
Corollary~\ref{cor:weighted} and the correction give
Theorem~\ref{thm:closed}. Proposition~\ref{prop:extraction} turns this
into the data (DL), and Proposition~\ref{prop:dl} gives a certificate,
proving Theorem~\ref{thm:universal} for $m\ge2$. Lemma~\ref{lem:padding}
covers $m=0,1$. Apply this to $A_0=\ZZ[C_1,\ldots,C_M]$ and transport to
$E$ (Section~\ref{sec:descent}); Lemma~\ref{lem:cert-conclusion}
concludes. \qed

\medskip
All choices in the proof are finite: lifts of finitely many
generators and coefficients, finitely many chart-wise lifts, finite
supports, and one mere partner at a time inside the noetherian
induction. The sequence $(z_r)$ is given by an explicit formula, so no
countable choice is needed. Excluded middle is not used; the only
decisions are on natural numbers, exponents, and membership in spans
in $D^v$ (strong discreteness).

\section{The formalisation}\label{sec:formal}

The proof is checked in Unit, the proof checker of
\url{https://github.com/coquand/unit}, a dependent type theory with
segments and hidden arguments. The design is by Coquand and the code
was written by Claude. All files are in \code{unit/sag/}.

\subsection{Foundations}
The development is parametrised by the segment \code{Theory} of
\code{sag/}: univalence, replacement (images) and surjections. Function
extensionality and propositional univalence are derived from these. No
new axiom is introduced. There is no choice principle: finite choice
(\code{FinCat.fchoice}) is proved by induction. There is no excluded
middle, countable choice, synthetic algebraic geometry axiom or
resizing. ``Merely'' is propositional truncation
(\code{Truncation.Trunc}), and rings are set-rings. Equality in $R$ is
never decided.

\subsection{The main theorem}
The final file is \code{zmt-general/FinalZMTGeneral.seg}:
\begin{verbatim}
module Theorem (R : CRing) (setR : isSet R.R)
 (I : IdealPred[R]) (propI : (r : R.R) -> isProp (I.In r))
 (Ea : Algebra over R) (n : Nat) (x : Fin n -> Ea.ring.R)
 (gen : (e : Ea.ring.R) -> || (H : Expr (Fin n)) * Id (eval x H) e ||)
 (fo : || QuotientFibre.Over[R,setR,I,propI,Ea].Finite ||) {
 Conclusion = (s : E) * (InIE (s - 1) *
     (ModuleFinite[R, B_s] * (InB_s s * isEquiv (phi : B_s[1/s] -> E[1/s]))))
 result : || Conclusion || }
\end{verbatim}
(abbreviated). Here $B_s=R[s,sx_1,\ldots,sx_n]$ is the image subalgebra,
with carrier a subtype of $E$. The map $\varphi$ is the canonical ring map
of localisations, and the conclusion states that this particular map is
an equivalence. Finiteness of $E/IE$ is \code{ModuleFinite} over the
quotient rings $R/I$ and $E/IE$ of \code{cffr/QuotientRing}. The same
file proves
\begin{verbatim}
universal : (n : Nat) -> CFNDescent.UniversalZMT n
\end{verbatim}
which is Theorem~\ref{thm:universal} at $A_0=\ZZ[C_1,\ldots,C_M]$: for
all residual data $\mathcal D_0$, merely a \code{ZMTCert}. Earlier
milestones had this principle as an explicit parameter; it is now
proved, and \code{result} has no parameter beyond those displayed.

\subsection{Correspondence with the proof}
Each directory is a separate checked development, linked to its
predecessors by relative symbolic links; earlier files are unchanged.

{\small\begin{longtable}{>{\raggedright\arraybackslash}p{0.3\textwidth}>{\raggedright\arraybackslash}p{0.63\textwidth}}
\toprule
Mathematics & Unit (directory / module) \\
\midrule
\endhead
Lemma~\ref{lem:cert-conclusion} & \code{zmt-descent}: \code{CertificateSubalgebra},
 \code{CertificateLocalisation}, \code{FiniteFibreLocalisation}\\
Lemma~\ref{lem:trick} & \code{exterior/DeterminantTrick},
 \code{zmt-descent/CharacteristicMonic}\\
Proposition~\ref{prop:dl} & \code{zmt-certificate-extract-general}:
 \code{GenericDegreeBounds} (dichotomy), \code{GenericFilteredSpan}
 (the module $V$), \code{GenericDegreeLoweringCertificate}\\
Lemma~\ref{lem:residual}, $E/IE$ finite & \code{zmt-descent}:
 \code{FiniteFibreCertificates.residual}, \code{QuotientFibre.bridge}\\
Universal data, transport & \code{zmt-descent}: \code{ResidualAlgebra},
 \code{UniversalResidualAlgebra}, \code{CoefficientBase} ($\ZZ[C]$ is CFN),
 \code{CertificateTransport}, \code{CFNDescent}\\
Lemma~\ref{lem:padding} & \code{zmt4-completion/ResidualPadding},
 \code{zmt-general/UniversalArityGlue}\\
Homogeneous equations & \code{zmt-sections-general}: \code{HomogeneousTermsGeneral},
 \code{GenericMonicSetup}\\
Chart rings, \v Cech cochains & \code{zmt-laurent-general} (actual localisations
 of the polynomial tower, Laurent coefficients, restrictions);
 \code{zmt-cech-general} (\code{GenericLaurent}, \code{GenericCube},
 \code{GenericCech})\\
$JC^p$, $I^rC^p$ & \code{zmt-rees-general}: \code{ImageQuotient},
 \code{GenericReesChart}; \code{FiniteIdealPowers} witness trees\\
Proposition~\ref{prop:sections} & \code{zmt-sections-general}:
 \code{GenericMonicChartRelation} \eqref{eq:G}, \code{GenericSectionData},
 \code{GenericSectionErrors}, \code{GenericSectionAssembly},
 \code{GenericFiniteFibreSections}; \code{zmt-assembly-general}\\
Rees kernel list & \code{zmt-rees-general}: \code{GenericReesRings},
 \code{GenericReesList}, \code{GenericReesExchange}
 (reusing \code{ReesFiniteList4} and \code{DegreeCor})\\
Lemma~\ref{lem:localised}, lifts $c_r$ & \code{GenericLocalizedReesMembership},
 \code{GenericLocalizedReesLift}, \code{GenericReesCochainMembership},
 \code{GenericReesPowerLift}\\
Theorem~\ref{thm:dependence}(a) & \code{zmt-cech-general} (contraction
 $dh+hd+\pi=\mathrm{id}$)\\
Theorem~\ref{thm:dependence}(b) & \code{zmt-syzygies-general}
 (\code{HomogeneousKernel}), \code{zmt-certificate-general}
 (\code{FinitePolynomialResolution}, \code{LocalizedSyzygies})\\
Theorem~\ref{thm:dependence}(c) & \code{zmt-certificate-general}:
 \code{CertificateExtension}, \code{FiniteZigzagCertificate},
 \code{CechZigzagSchedule}, \code{ScalarResolutionDependence}\\
Theorem~\ref{thm:dependence}(d) & \code{zmt-lifting/CertifiedDependence}\\
Corollary~\ref{cor:weighted} & \code{zmt-weighted-general}:
 \code{GenericWeightedDependence}, \code{EvaluatedWeightedDependence}\\
Correction, Theorem~\ref{thm:closed} & \code{zmt-rees-general}:
 \code{GenericReesFactor}, \code{GenericReesCorrection},
 \code{GenericReesFromSections}; \code{zmt-correction-general}\\
Proposition~\ref{prop:extraction} & \code{zmt-extract-exponents-general}
 (\code{ExtractionBounds}), \code{zmt-extract-general}
 (\code{PrescribedPolynomialization}, \code{GenericClearedExtraction}),
 \code{zmt-general/GenericSectionExtraction}\\
Assembly & \code{zmt-general/FinalZMTGeneral}\\
\bottomrule
\end{longtable}}

\subsection{Representation choices}
Some formal choices differ from the text, without changing the
mathematics.
\begin{itemize}
\item Exponents are functions $\mathrm{Fin}\,n\to\mathbb N$ or
$\mathrm{Fin}\,n\to\ZZ$. Subsets of coordinates are structural bitsets.
The newest coordinate is $Z$ and comes first in the orientation, so
$d(a,b)=(da,a-db)$ uniformly in $n$, with no sign tables.
\item Chart elements are finitely supported Laurent coefficient
functions with a support condition. The native localisation of the
polynomial tower is compared with them by coefficient extensionality.
\item Ideal powers $I^r$ are explicit finite witness trees; $J$-membership
is an image of the action of the original $\mathbf P_q$.
There are no quotient modules and no tensor products.
\item The polynomial ring is an iterated tower $A[X_1]\cdots[X_n]$.
Homogeneous elements are handled through finite term lists, and
$\deh$ is evaluation of the tower. Monic polynomials are data with a
literal leading coefficient $1$.
\item The dependence theorem is proved for $\nu=\mathit{extra}+3$
coordinates and the sections for $m+1$ coordinates; one transport along
$\mathit{extra}+3=(\mathit{extra}+2)+1$ (\code{GenericCorrectionTransport})
identifies the two interfaces.
\end{itemize}

\subsection{Validation}
On 30 October 2026, a fresh run of the checker on
\code{zmt-general/All.seg}, which loads every development above,
checked $18\,779$ declarations with $0$ failures in $13.1$\,s. The
checking server reported $359$ modules, $0$ failed. The opening audit
(\code{--opens}) confirms, for every declaration owned by the final
directory, that it unfolds only the listed definitions. Each owned
declaration opens at most $35$ definitions and takes at most $0.02$\,s.
The intermediate directories carry their own audits; for example
\code{zmt-rees-general} owns $1121$ declarations and
\code{zmt-certificate-general} owns $1132$. All commands were bounded to
$55$\,s.

The regressions run the final theorem with no assumed principle:
\begin{itemize}
\item $R=E=\ZZ/4$, $I=(2)$, every generator equal to $2$, for
$n=0,\ldots,5$, so the base is not reduced;
\item the zero fibre $I=R$, for $n=0$ and $n=4$;
\item extraction with no affine variables, and extraction from the
constructed correction for arbitrary residual data in four variables;
\item certificates over $\ZZ$ for $m=0,4$, and over $\ZZ/4$ with
nonzero nilpotent generators.
\end{itemize}
Deliberate mutations are rejected. The final conclusion with $s+1$ in
place of $s-1$ fails, and so does the zero map in place of the canonical
localisation map. Earlier tracks reject a reversed edge orientation, a
wrong homogeneous degree, the bound $d=D$ in place of $D+\sum a_i$, the
correction $c+ap$ in place of $c-ap$, and a sign error in a coefficient
degree shift.

\subsection{Scope}
The checked statement is Theorem~\ref{thm:main} for every finite
number of generators. Its hypotheses are an ideal with propositional
membership and the quotient-ring model of $E/IE$ in \code{cffr}. A caller
with a different quotient construction needs an isomorphism to this
one. The inverse of $\varphi$ is not packaged separately as a ring map;
it is one, since $\varphi$ is a bijective ring map. An earlier,
independent development (\code{zmt4-completion}) proves the case of at
most three generators with four fixed homogeneous coordinates; the
general proof uses only its dimension-independent helpers, such as the
padding lemma.

\begin{thebibliography}{9}
\bibitem{ACL} M.~Alonso, T.~Coquand and H.~Lombardi,
\emph{Revisiting Zariski Main Theorem from a constructive point of
view}, J. Algebra 406 (2014), 46--68.
\bibitem{CP} T.~Coquand and H.~Persson, \emph{Gr\"obner bases in type
theory}, Types for Proofs and Programs, LNCS 1657 (1999), 33--46.
\bibitem{Local} T.~Coquand, \emph{Local structure of \'etale algebras},
manuscript (\code{zmt2-revised.pdf}).
\bibitem{Stacks} The Stacks Project, Section 30.20, \emph{The theorem
on formal functions}.
\bibitem{Notes} Internal notes in \code{unit/doc/}:
\code{zmt-formal-cohomology.tex} (the plan of this proof) and the
\code{zmt-*-handoff.txt} files.
\end{thebibliography}
\end{document}
